\chapter{Gain Structure}
\label{chap:gain-structure}

Gain structure is about how large the signal should be at each stage. It is not
the same as the fader: gain sets the level entering the next stage, while the
fader sets how loud this track is relative to the others. Clean gain comes first;
the fader only makes sense after that.

\section{Level, gain, fader}

\begin{itemize}
  \item \term{Level}: the size of a signal at a given moment, usually expressed in
        dBFS or dB.
  \item \term{Gain}: amplification or attenuation applied by a device. Turning up
        gain makes the signal stronger, and also makes noise and distortion louder.
  \item \term{Fader}: the relative volume control on a channel strip. It changes
        this track's share of the bus, not ``fixes distortion''.
\end{itemize}

Treating the fader as a universal volume key is the most common mixing mistake:
when a track is too quiet, push the fader to +6 dB and discover that the device
chain was already clipping; when a track is too loud, pull the fader to −30 dB
and lose all the dynamics. The correct order is to put the device chain at a
sensible level first, then use the fader for relative balance.

\begin{center}
\resizebox{\textwidth}{!}{\begin{tikzpicture}[
  node distance=8mm and 6mm,
  box/.style={draw, rounded corners=2pt, minimum width=26mm, minimum height=8mm, align=center, font=\small},
  arrow/.style={-{Latex[length=2mm]}, thick}
]
  \node[box] (source) {Source / clip};
  \node[box, right=of source] (devices) {Device-chain gain\\EQ / dynamics / space};
  \node[box, right=of devices] (fader) {Fader\\relative volume};
  \node[box, right=of fader] (bus) {Bus\\summing};
  \node[box, right=of bus] (master) {Master\\monitor / export};

  \draw[arrow] (source) -- (devices);
  \draw[arrow] (devices) -- (fader);
  \draw[arrow] (fader) -- (bus);
  \draw[arrow] (bus) -- (master);
\end{tikzpicture}}
\end{center}

``Device-chain gain'' in the diagram is the working level inside the devices,
not the mixer fader. Yeban currently exposes no editable input-gain knob on the
channel strip; gain inside the device chain comes from each effect's own
parameters, covered in Chapters 3 and 4.

\section{Readings on Yeban's channel strip}

Each channel strip in the mixer console (X / F8) has the following controls and
readings. Ranges and defaults come from \texttt{TrackV3} and the fader mapping
code; the interface shows a lossy one-decimal display.

\begin{table}[htbp]
\centering
\small
\begin{tabular}{@{}p{0.20\linewidth}p{0.27\linewidth}p{0.44\linewidth}@{}}
\hline
Control/readout & Range and default & Behaviour \\
\hline
Track fader \texttt{volume\_db} & −60.0 to +6.0 dB; default 0.0 dB & Linear fader travel; 0 dB sits at about 90.9\% of travel; double-click resets to 0 dB \\
Pan \texttt{pan} & −1.0 to +1.0; default 0.0 & Drag left/right; double-click centres \\
Mute / solo & Boolean; default false & Audition tools; they do not delete content \\
Track meter & Normalised 0.0–1.0 & Comes from the engine \texttt{MeterFrame}; above 0.8 turns to the warning colour \\
Master meters & Left/right, each 0.0–1.0 & Mirrored from the master frame at 60 Hz \\
Monitoring loudness & Integrated/short-term LUFS, peak dBTP & Shown in the control room; labelled EBU R128 \\
\texttt{IN: 0.0dB} / \texttt{Ø} & Shown when the mixer is maximised & Currently a static display; no editable input gain or polarity-invert action is wired \\
\hline
\end{tabular}
\caption{Yeban channel-strip controls and readings}
\label{tab:channel-strip-readings}
\end{table}

\screenshot{mixer-channel-strip}{The mixer console is maximised and one channel strip is fully visible: name, IN/Ø, fader dB, pan, and meter.}

Fader position and dB are not an exponential relationship; the travel maps
linearly to −60.0 to +6.0 dB. Moving the fader slightly down from 0 dB changes dB
by very little; only near the bottom of travel does it reach −60 dB. Read the
number, not just the cap height.

The track meter is a normalised value, not a dBFS number. Use it to see whether a
signal exists and whether it is running hot; read precise peak and loudness in the
control room. Above 0.8 the meter changes to the warning colour, meaning ``near
the top'', not ``already clipping''. If it stays in the warning zone, fix it
before export.

\section{Three principles of gain structure}

\begin{enumerate}
  \item \term{Gain first, fader second}: set a sensible working level in the device
        chain, then balance with the fader. If the chain already clips, lowering the
        fader only makes the clipped waveform smaller; it does not repair it.
  \item \term{Leave headroom at every stage}: do not push everything to 0 dBFS.
        Leave space for dynamics, parallel summing, and bus processing; during
        mixing, peak headroom usually matters more than ``looking louder''.
  \item \term{Amplify only when needed}: louder is not automatically better. When a
        track needs to be louder, first check arrangement, timbre, and balance;
        blind gain boosts noise, distortion, and masking together.
\end{enumerate}

``0 dBFS'' is digital full scale, not a standard volume. Yeban's fader 0 dB is
only a project reference point; it does not mean the signal is at 0 dBFS. Actual
level is determined by the source, the device chain, and the fader together.

\section{Stage-by-stage check}

\begin{enumerate}
  \item Solo the current track and confirm there is signal; the meter should not be
        at zero. If there is no signal, check the clip, mute/solo state, and device
        bypass first; do not reach for the fader.
  \item Check whether the track meter stays above 0.8. If it does, go back to the
        device chain and inspect gain and effect parameters first, then bring the
        fader to a sensible position.
  \item Balance with faders: start with the lead vocal, drums, and bass skeleton,
        then place other tracks around them. Do not push every track from first to
        last.
  \item Check the master bus and control room: whether the master meter is pinned
        near the top, whether integrated/short-term LUFS is in the same range as the
        reference, and whether peak dBTP has headroom.
  \item Before export, check mute/solo, the reference track, and bypassed devices
        once more. Any one of them out of place can make the export differ from what
        you heard in the project.
\end{enumerate}

\section{Common mistakes}

\begin{itemize}
  \item \term{Using the fader to fix clipping}: the fader can only change the size
        of an already-clipped signal; it cannot restore the waveform.
  \item \term{Pushing every fader up}: if every track goes +6 dB, the master bus
        only clips sooner; relative balance does not require an overall boost.
  \item \term{Confusing monitoring volume with mix volume}: DIM/MONO/SIDE change
        only how you monitor. The master fader is project state; do not conflate them.
  \item \term{Treating IN:0.0dB as an input-gain knob}: it is currently only a
        static readout in the maximised mixer, with no editable action. To change
        device gain, edit the relevant effect parameters in the device chain.
\end{itemize}

\section{Exercise}

\begin{enumerate}
  \item Choose a track with level variation. Solo it and confirm the meter is not at zero; then check ``device-chain gain $\to$ fader $\to$ bus'' in order.
  \item Pull the fader from 0 dB to $-12$ dB, then restore the same loudness in the device chain; compare transients and noise floor.
  \item Note one section where the meter enters the 0.8 warning zone; check the device chain first, then use the fader for relative balance. Do not fix clipping with the fader.
\end{enumerate}

\section{Help-topic mapping for this chapter}

\begin{table}[htbp]
\centering
\small
\begin{tabular}{@{}p{0.34\linewidth}p{0.52\linewidth}@{}}
\hline
Section & In-app help anchor \\
\hline
Level, gain, fader & \texttt{gain-level-fader} \\
Readings on Yeban's channel strip & \texttt{gain-channel-strip} \\
Three principles of gain structure & \texttt{gain-structure-principles} \\
Stage-by-stage check & \texttt{gain-check-workflow} \\
Common mistakes & \texttt{gain-common-mistakes} \\
\hline
\end{tabular}
\caption{Help-topic mapping for Chapter 2}
\end{table}
